\documentclass[12pt]{article}
\usepackage[utf8]{inputenc}
\usepackage[T1]{fontenc}
\usepackage{hyperref}
\title{Lemnian Language Resources for K-12 Curriculum Development}
\author{Compiled by Claude Sonnet 5.5 (claude-sonnet-5-5) for Andr\'as Kornai}
\date{}
\begin{document}
\maketitle

\section{Introduction}
Lemnian is an extinct language attested on the island of Lemnos in the second half of the 6th century BC. It is known mainly from one funerary stele (the Lemnos stele, found in 1885 near Kaminia), a few pottery fragments, and a further inscription reported from Hephaistia in 2009. Most scholars place it in a Tyrsenian group with Etruscan and Raetic. It was written in a Greek-derived alphabet and was replaced by Attic Greek after Athens took the island. The corpus is tiny, there are no native speakers, and no curriculum materials were found in my searches. A K-12 programme using the language itself is not realistic beyond a few attested words and phrases, so any such programme would have to be a heritage or history unit on Lemnos.

\section{Primary Educational Resources}
No primers, textbooks or school programmes for Lemnian were found in search results.

\section{Linguistic and Reference Materials}
\begin{itemize}
\item Wikipedia, ``Lemnian language'': an overview of the stele, the pottery evidence, classification and script. \url{https://en.wikipedia.org/wiki/Lemnian_language}
\item Verbix documentation page on Lemnian. I did not check its content. \url{https://docs.verbix.com/Languages/Lemnian}
\end{itemize}

\section{Religious and Community Resources}
No community organisations or religious resources for Lemnian were found. It has no living community.

\section{Digital and Audio Resources}
\begin{itemize}
\item Greek Reporter, an article on the discovery of the Lemnos stele and its retention in Greece (2025). It is useful as a narrative hook for pupils. \url{https://greekreporter.com/2025/11/19/lemnian-stele-story-discovery-intrigue-patriotism/}
\end{itemize}
No audio or software resources were found.

\section{Curriculum Development Context}
There are only a handful of short inscriptions, and their interpretation is uncertain. A curriculum could cover the Lemnos stele as a historical and archaeological topic, the Greek-derived script, and the question of Tyrsenian relationships. Material from specialist epigraphic sources would need to be located separately. The language cannot be taught as a spoken language.

\end{document}